\begin{resumo}
Assim, o presente pré-projeto situa-se na interface entre Administração Pública, gestão de despesas institucionais e Ciência de Dados aplicada à tomada de decisão. A pesquisa parte da premissa de que a gestão de passagens aéreas pode ser aprimorada por meio de uma abordagem analítica de apoio ao planejamento, com componente preditivo exploratório, capaz de combinar registros internos, dados públicos e monitoramento prospectivo de preços. Ao utilizar dados da [instituição-caso anonimizada] como base empírica e estruturar uma coleta controlada de cotações, busca-se desenvolver um painel ou protótipo analítico que indique padrões de variação tarifária, faixas de antecedência e níveis de incerteza por rota, com potencial de replicação por órgãos e entidades públicas que disponham de dados compatíveis ou possam implementar rotinas semelhantes de monitoramento.
A comparação considera que compras realizadas não revelam preços alternativos da mesma passagem; por isso o monitoramento prospectivo mantém rota, data de viagem e condições de busca comparáveis. A economia pretendida permanece hipótese a avaliar, sem resultado empírico declarado neste esqueleto.
\end{resumo}
